\subsection*{Question 2.8 - Relation between attitude and horizontal channels}

\paragraph{Relation between the transfer functions.}
Dividing the horizontal transfer functions of Question 2.7 by the attitude ones, the inertia and the input cancel and only a double integrator and the gravity remain:
\begin{equation*}
G_y(s) = \frac{g}{s^2}\, G_\phi(s), \qquad
G_x(s) = -\frac{g}{s^2}\, G_\theta(s).
\end{equation*}
In words, each horizontal channel is the corresponding attitude channel cascaded with a gain $\pm g$ and two integrators.

\paragraph{Physical interpretation.}
The rotors can only produce a force along $-z_B$.
So the vehicle has no direct way of pushing itself sideways, and to accelerate horizontally it must first tilt, so that a component of the thrust points in the desired direction.
When the thrust balances the weight ($T \approx m g$), a small roll angle $\phi$ tilts the thrust vector towards $+y_B$ and produces a horizontal acceleration of $g\phi$ along $y_B$.
A small pitch angle $\theta$ tilts the thrust towards $-x_B$ and produces an acceleration $-g\theta$ along $x_B$, which explains the minus sign in $G_x$.
The position is the double integral of the acceleration, which is the factor $g/s^2$ above.
This means that the horizontal position can only be controlled indirectly through the attitude: the angle has to build up first (its own double integrator), and then the position reacts.
This is the reason for the cascade structure of the controller, with a fast inner attitude loop and a slower outer position loop.
It also shows that the horizontal channels are harder to control than the vertical one, as they have relative degree four (four integrators) instead of two.

\paragraph{Effect of the yaw equilibrium $\psi_0$.}
The matrices $A$ and $B$ of Question 2.6 do not depend on $\psi_0$.
This is because the error $\tilde p = R^T(\lambda)(p - p_0)$ is expressed in body-aligned axes, and the dependence on the heading is removed by the rotation $R^T$.
Consequently the transfer functions $G_\phi, G_\theta, G_\psi, G_x, G_y, G_z$, and the relations above, are exactly the same for every choice of $\psi_0$.

What does depend on $\psi_0$ is the meaning of $\tilde x$ and $\tilde y$ in the inertial frame.
To first order, $p - p_0 = R(\lambda)\,\tilde p \approx R_z(\psi_0)\,\tilde p$, so
\begin{equation*}
\begin{bmatrix} \Delta p_N \\ \Delta p_E \end{bmatrix} \approx
\begin{bmatrix} \cos\psi_0 & -\sin\psi_0 \\ \sin\psi_0 & \cos\psi_0 \end{bmatrix}
\begin{bmatrix} \tilde x \\ \tilde y \end{bmatrix},
\qquad \Delta p_D = \tilde z .
\end{equation*}
For $\psi_0 = 0$ the body-aligned and inertial axes coincide, so a roll moves the vehicle along the inertial $y$ axis and a pitch along the inertial $x$ axis.
For $\psi_0 \neq 0$ the pairing $G_\phi \leftrightarrow G_y$, $G_\theta \leftrightarrow G_x$ still holds, but only along the heading-aligned axes: a roll command moves the vehicle along the direction of $y_B$, which is rotated by $\psi_0$ with respect to the inertial $y$ axis.
The vertical channel is not affected, since $z$ is the rotation axis of the yaw.
In practice, an outer position loop working in the inertial frame must rotate the desired horizontal accelerations by $\psi_0$ before turning them into roll and pitch references.
